\chaptertitle{第四章\quad 数据的波动}

第三章我们学了平均数、中位数、众数——它们告诉我们数据"聚集在哪里"。但光知道聚集在哪里还不够。想象两个班的数学平均分都是80分——一班每个同学都在75~85分之间，二班有的考30分、有的考100分。这两个班的成绩分布截然不同。这一章就来研究数据"分散"的程度——波动性。

\sectiontitle{4.1\quad 为什么要描述波动？}

先看两个具体例子来建立直觉。

\begin{example}
甲乙两名射击运动员各打了5发子弹，成绩（环数）：\\
甲：8, 8, 8, 8, 8 \qquad 乙：6, 7, 8, 9, 10\\
平均成绩相同吗？谁的发挥更稳定？
\end{example}
\begin{solution}
甲的平均数$=(8+8+8+8+8)\div5=8$环。\\
乙的平均数$=(6+7+8+9+10)\div5=8$环。\\
两人平均成绩相同——但甲每枪都是8环，发挥极其稳定；乙忽高忽低，不稳定。仅靠平均数无法区分这两种情况，需要一个新的量来描述"波动大小"。
\end{solution}

\begin{center}
\begin{tikzpicture}[scale=0.8]
\draw[->,thick] (0,0) -- (7,0) node[right]{成绩（环）};
\draw[->,thick] (0,0) -- (0,2.5) node[above]{};
\foreach \x in {6,7,8,9,10} {
  \draw (\x-6+0.5,0.1) -- (\x-6+0.5,-0.1) node[below]{$\x$};
}
% 甲 - all at 8
\foreach \x in {0.8,1.5,2.2,2.9,3.6} {
  \filldraw[mainblue] (\x,1.5) circle (3pt);
}
\node[mainblue] at (2.2,2){甲：全部集中在8环};
% 乙 - spread out
\foreach \x/\v in {0.5/6,1.2/7,2.2/8,3.2/9,3.9/10} {
  \filldraw[exred] (\x,0.5) circle (3pt);
}
\node[exred] at (2.2,0.2){乙：散布在6~10环};
\draw[dashed] (0,1) -- (7,1);
\node at (7.3,1){平均数8环};
\node at (3.2,-1.2){\small 图4-1\quad 甲和乙的成绩分布对比};
\end{tikzpicture}
\end{center}

\begin{example}
两名同学近5次数学测验：小明——70, 80, 90, 80, 80；小红——50, 90, 100, 70, 90。平均分都是80分。谁的波动更大？
\end{example}
\begin{solution}
从数据直接看：小明的成绩在70~90之间波动；小红在50~100之间。显然小红波动更大。需要量化这种差异。
\end{solution}

\sectiontitle{4.2\quad 方差——偏离平均的程度}

怎么量化"波动大小"？一个自然的想法是：看每个数据离平均数有多远，然后取"平均偏离距离"。

先看甲的5发成绩（全为8，平均数=8）：
\begin{itemize}
\item 每个数据与平均数的差（\textbf{离差}）：$8-8=0$（5次）
\item 离差的平均：$(0+0+0+0+0)\div5=0$
\end{itemize}

再看乙的成绩（平均数=8）：
\begin{itemize}
\item 离差：$6-8=-2$，$7-8=-1$，$8-8=0$，$9-8=1$，$10-8=2$
\item 离差的平均：$(-2-1+0+1+2)\div5=0$
\end{itemize}

离差的平均竟然也是0！因为正负离差互相抵消了。解决方法是把离差\textbf{平方}（让所有值变成非负数）再取平均——这就是\textbf{方差}。

甲的方差：
\[
s^2 = \frac{0^2+0^2+0^2+0^2+0^2}{5}=0
\]

乙的方差：
\[
s^2 = \frac{(-2)^2+(-1)^2+0^2+1^2+2^2}{5} = \frac{4+1+0+1+4}{5}=2
\]

甲的方差为0（完全没波动），乙的方差为2（有波动）。方差越大，数据越分散。

\knowbox{
\textbf{方差}（$s^2$）：\\
$s^2=\dfrac{(x_1-\bar x)^2+(x_2-\bar x)^2+\cdots+(x_n-\bar x)^2}{n}$。\\
方差越大 $\rightarrow$ 数据越分散（波动大）；\\
方差为0 $\rightarrow$ 所有数据完全相同。}

\begin{example}
计算数据$3,5,7,9,11$的方差。
\end{example}
\begin{solution}
平均数$\bar x=(3+5+7+9+11)\div5=7$。\\
$s^2=\dfrac{(3-7)^2+(5-7)^2+(7-7)^2+(9-7)^2+(11-7)^2}{5}
=\dfrac{16+4+0+4+16}{5}=8$。
\end{solution}

\begin{example}
两组数据：A组$5,5,5,5,5$；B组$1,3,5,7,9$。分别求方差，哪组更稳定？
\end{example}
\begin{solution}
A组平均数$=5$，方差$=\dfrac{0+0+0+0+0}{5}=0$（完全稳定）。\\
B组平均数$=5$，方差$=\dfrac{(1-5)^2+(3-5)^2+(5-5)^2+(7-5)^2+(9-5)^2}{5}
=\dfrac{16+4+0+4+16}{5}=8$。\\
A组更稳定。事实上，A组的数据全部相等，没有任何波动。
\end{solution}

\begin{exercise}
\item 计算数据$2,4,6,8,10$的方差。
\item 两组数据：P组$80,90,80,90,80$；Q组$70,80,90,80,90$。哪组波动更大？用方差说明。
\item 一组数据的方差是$0$，这意味着什么？请构造一组$5$个数的数据使其方差为$0$。
\end{exercise}

\sectiontitle{4.3\quad 标准差——回到原来的单位}

方差有一个小问题：单位是原始数据单位的平方。比如身高数据的单位是cm，方差的单位就是$cm^2$——这让人不太舒服。为了把单位还原回来，我们对方差取平方根，得到\textbf{标准差}。

\knowbox{
\textbf{标准差}：$s=\sqrt{s^2}$（方差的算术平方根）。\\
标准差的单位与原始数据相同，更容易解释。\\
例如：平均身高$170$cm，标准差$10$cm——意思是大多数人的身高在$170\pm10$cm范围内。}

\begin{example}
计算乙组数据（$6,7,8,9,10$）的标准差。
\end{example}
\begin{solution}
方差$s^2=2$，标准差$s=\sqrt2\approx1.41$环。\\
可以这样说：乙的每发成绩平均偏离平均成绩约$1.41$环。
\end{solution}

\begin{example}
某班学生身高方差$100$（$cm^2$），求标准差并解释其含义。
\end{example}
\begin{solution}
标准差$=\sqrt{100}=10$cm。可以理解为：大部分学生的身高在平均身高的$\pm10$cm范围内波动。
\end{solution}

\begin{exercise}
\item 数据$10,12,14,16,18$的方差和标准差各是多少？
\item 两组数据：A的标准差$s_A=5$，B的标准差$s_B=10$。哪组数据更分散？为什么？
\item 一组数据的标准差为$0$，这组数据有什么特征？
\end{exercise}

\sectiontitle{4.4\quad 用方差做比较与决策}

方差和标准差在实际中广泛用于比较稳定性和一致性——从产品质量到运动员表现。

\begin{example}
两家供应商提供同型号电池寿命（小时）：\\
A厂：98, 101, 99, 102, 100 \qquad B厂：90, 105, 95, 110, 100\\
如果要求电池寿命\textbf{稳定}（即每节电池的寿命差异不能太大），你会选哪家？
\end{example}
\begin{solution}
A厂平均数$=100$，方差$=\dfrac{(-2)^2+1^2+(-1)^2+2^2+0^2}{5}=2$。\\
B厂平均数$=100$，方差$=\dfrac{(-10)^2+5^2+(-5)^2+10^2+0^2}{5}=50$。\\
A厂方差远小于B厂，说明A厂电池寿命更稳定，应该选A厂。虽然两厂的平均寿命相同，但B厂的电池有的只能用90小时、有的能用110小时——对需要稳定供电的设备来说，这是不可接受的。
\end{solution}

\begin{example}
两名射箭运动员的成绩（环数）：\\
甲：8, 9, 8, 9, 8 \qquad 乙：7, 10, 8, 9, 8\\
(1) 谁的平均成绩更高？(2) 谁的成绩更稳定？
\end{example}
\begin{solution}
(1) 甲平均$=(8+9+8+9+8)\div5=8.4$；乙平均$=(7+10+8+9+8)\div5=8.4$。平均成绩相同。\\
(2) 甲的方差$=\dfrac{(-0.4)^2+0.6^2+(-0.4)^2+0.6^2+(-0.4)^2}{5}=0.24$。\\
乙的方差$=\dfrac{(-1.4)^2+1.6^2+(-0.4)^2+0.6^2+(-0.4)^2}{5}=1.04$。\\
甲更稳定（方差更小）。比赛中，稳定的选手往往比时好时差的选手更可靠。
\end{solution}

\begin{center}
\begin{tikzpicture}[scale=0.7]
\draw[->,thick] (0,0) -- (7,0) node[right]{环数};
\draw[->,thick] (0,0) -- (0,2.5) node[above]{};
\foreach \x in {7,8,9,10} {
  \draw (\x-6+0.5,0.1) -- (\x-6+0.5,-0.1) node[below]{$\x$};
}
\draw[dashed] (0,1.1) -- (7,1.1);
\node at (7.3,1.1){平均8.4环};
% 甲
\foreach \x in {1.0,1.8,2.6,3.4,4.2} {
  \filldraw[mainblue] (\x,1.8) circle (3pt);
}
\node[mainblue] at (2.6,2.3){甲（集中）};
% 乙
\foreach \x/\v in {0.7/7,1.5/8,2.6/8,3.7/9,4.5/10} {
  \filldraw[exred] (\x,0.4) circle (3pt);
}
\node[exred] at (2.6,-0.3){乙（分散）};
\node at (3.5,-1.2){\small 图4-2\quad 甲与乙的成绩离散程度对比};
\end{tikzpicture}
\end{center}

\begin{example}
某车间两名工人加工零件，要求直径$10$mm。各加工$5$个零件的直径（mm）：\\
王师傅：10.0, 10.1, 9.9, 10.0, 10.0 \qquad 李师傅：10.2, 9.8, 10.1, 9.9, 10.0\\
谁加工的零件更符合要求（更稳定）？
\end{example}
\begin{solution}
王师傅平均$=10.0$，方差$=\dfrac{0+0.01+0.01+0+0}{5}=0.004$。\\
李师傅平均$=10.0$，方差$=\dfrac{0.04+0.04+0.01+0.01+0}{5}=0.02$。\\
王师傅的方差更小——他的零件直径更稳定，质量更可靠。
\end{solution}

\begin{exercise}
\item 两个班级期末数学平均分都是78分，一班标准差6分，二班15分。从"教学质量是否整齐"的角度看，哪个班更好？
\item 某股票过去5天的日收益率（\%）：$-2, +5, -1, +3, -5$。另一支股票：$+1, +2, 0, -1, +3$。哪支股票波动更大？你会建议风险承受能力低的人选哪支？
\item 数据$a-3,a-1,a,a+1,a+3$的方差是多少？提示：先求平均数。
\end{exercise}

\sectiontitle{本章小结}

平均数和方差是描述一组数据的两个方面：平均数告诉我们"数据在哪"，方差告诉我们"数据有多散"。方差是每个数据与平均数的差的平方的平均——它通过平方避免了正负抵消。标准差是方差的平方根，把单位还原回来，更便于解释。在实际中——产品质量控制、运动员选拔、投资风险评估——方差和标准差是不可或缺的工具。一组数据的完整描述应该是"均值$\pm$标准差"，而不是孤零零的一个平均数。

\sectiontitle{课后练习}

\subsection*{A组\quad 基础巩固}

\begin{exercise}
\item 计算数据$4,7,10,13,16$的方差和标准差。
\item 数据组的方差是$9$，标准差是多少？如果每个数据都加上$5$，新数据的方差是多少？
\item 两组数据：A组$2,5,8,11,14$；B组$6,7,8,9,10$。\\
(a) 分别求平均数和方差；(b) 哪组更集中？
\item 甲、乙各投篮$10$次（1=命中，0=未中）：\\
甲：1,1,1,1,1,0,0,0,1,1 \qquad 乙：1,0,1,0,1,0,1,0,1,1\\
求两人的命中率和命中次数的方差。谁的发挥更稳定？
\item 一组数据的标准差是$4$，如果每个数据都乘以$2$，新数据的标准差是多少？
\end{exercise}

\subsection*{B组\quad 挑战提升}

\begin{exercise}
\item 验证：数据$3,5,7,9,11$的方差是$8$。如果每个数据都加上同一个常数$a$，方差不变；如果都乘以$k$，方差变为原来的$k^2$倍。试解释为什么。
\item 两组数据$x_1,\dots,x_n$和$y_1,\dots,y_n$，平均数分别为$\bar x,\bar y$，方差分别为$s_x^2,s_y^2$。合并后的总方差不一定等于$(s_x^2+s_y^2)/2$。举例说明。
\item 某比赛评委打分：$9.5,9.2,9.8,9.0,9.6,9.4$。去掉一个最高分和一个最低分后，剩下分数的方差是变大还是变小了？为什么体育比赛常用"去掉最高最低分"的做法？
\item 数据$x_1,\dots,x_n$的方差是$s^2$。令$y_i=x_i-b$，则$y_i$的方差是多少？令$z_i=kx_i$呢？请用公式推导。
\item 甲乙两人进行投篮比赛，每人投$100$次，甲命中$80$次，乙命中$75$次。但在$10$次一组的统计中，甲的每次命中率在$0.7$~$0.9$之间波动，乙在$0.73$~$0.77$之间波动。虽然甲的平均命中率更高，但为什么在关键比赛中教练可能会更信任乙？
\end{exercise}

\newpage
